\chapter*{Appendix D: Key Risk Assessment Frameworks and Standards}
\label{chap:Appendix D: Key Risk Assessment Frameworks and Standards}

%---------------------------- SECTION ----------------------------
\section*{D.1 General Risk Management}

\begin{table}[H]
  \centering
  \renewcommand{\arraystretch}{1.5}
  \begin{tabularx}{\textwidth}{ | >{\raggedright\arraybackslash}p{0.2\textwidth} 
								| >{\raggedright\arraybackslash}p{0.2\textwidth} 
								| >{\raggedright\arraybackslash}p{0.4\textwidth} 
                                | >{\raggedright\arraybackslash}X | }
    \hline
    \textbf{Framework / Standard} & \textbf{Published by} & \textbf{Key Purpose} & \textbf{Primary Audience} \\
    \hline
    \textbf{ISO 31000:2018} & ISO & International standard for risk management — principles, framework, and process & All organisations\\
    \hline
    \textbf{ISO 31010:2019} & ISO & Risk assessment techniques — catalogue and guidance on 31 techniques & Risk practitioners\\
    \hline
    \textbf{COSO ERM Framework (2017)} & COSO & Enterprise risk management — integrating with strategy and performance & Corporate governance\\
    \hline
    \textbf{Orange Book (2023)} & HM Treasury & Risk management guidance for UK central government departments & Public sector\\
    \hline
    \textbf{IRM Risk Management Standard} & Institute of Risk Management & Foundational risk management standard & Risk professionals\\
    \hline
  \end{tabularx}
  \caption{General Risk Management}
\end{table}

%---------------------------- SECTION ----------------------------
\section*{D.2 Health, Safety, and Environment}

\begin{table}[H]
  \centering
  \renewcommand{\arraystretch}{1.5}
  \begin{tabularx}{\textwidth}{ | >{\raggedright\arraybackslash}p{0.4\textwidth} 
								| >{\raggedright\arraybackslash}p{0.2\textwidth} 
                                | >{\raggedright\arraybackslash}X | }
    \hline
    \textbf{Framework / Standard} & \textbf{Published by} & \textbf{Key Purpose} \\
    \hline
    \textbf{HSE Five Steps to Risk Assessment} & HSE & Practical guidance on the risk assessment process for workplace hazards\\
    \hline
    \textbf{HSE Management Standards for Work-Related Stress} & HSE & Framework for assessing and managing psychosocial risk factors\\
    \hline
    \textbf{ISO 45001:2018} & ISO & International standard for occupational health and safety management systems\\
    \hline
    \textbf{ISO 14001:2015} & ISO & International standard for environmental management systems\\
    \hline
    \textbf{COSHH Essentials} & HSE & Guidance on assessing and controlling risks from hazardous substances\\
    \hline
  \end{tabularx}
  \caption{Health, Safety, and Environment}
\end{table}

%---------------------------- SECTION ----------------------------
\section*{D.3 Cybersecurity and Information Security}

\begin{table}[H]
  \centering
  \renewcommand{\arraystretch}{1.5}
  \begin{tabularx}{\textwidth}{ | >{\raggedright\arraybackslash}p{0.3\textwidth} 
								| >{\raggedright\arraybackslash}p{0.2\textwidth} 
                                | >{\raggedright\arraybackslash}X | }
    \hline
    \textbf{Framework / Standard} & \textbf{Published by} & \textbf{Key Purpose} \\
    \hline
    \textbf{NIST Cybersecurity Framework (CSF) 2.0} & NIST & Voluntary framework for managing cybersecurity risk — Identify, Protect, Detect, Respond, Recover\\
    \hline
    \textbf{ISO 27001:2022} & ISO & International standard for information security management systems\\
    \hline
    \textbf{ISO 27005:2022} & ISO & Guidance on information security risk management\\
    \hline
    \textbf{SOC 2} & AICPA & Framework for assessing controls relevant to security, availability, processing integrity, confidentiality, and privacy\\
    \hline
    \textbf{Cyber Essentials} & NCSC & UK government-backed scheme defining a baseline of cybersecurity controls\\
    \hline
    \textbf{NCSC 10 Steps to Cyber Security} & NCSC & Guidance on the most effective steps organisations can take to improve their cybersecurity\\
    \hline
  \end{tabularx}
  \caption{Cybersecurity and Information Security}
\end{table}

%---------------------------- SECTION ----------------------------
\section*{D.4 Financial Services}

\begin{table}[H]
  \centering
  \renewcommand{\arraystretch}{1.5}
  \begin{tabularx}{\textwidth}{ | >{\raggedright\arraybackslash}p{0.35\textwidth} 
								| >{\raggedright\arraybackslash}p{0.25\textwidth} 
                                | >{\raggedright\arraybackslash}X | }
    \hline
    \textbf{Framework / Standard} & \textbf{Published by} & \textbf{Key Purpose} \\
    \hline
    \textbf{Basel Framework} & BCBS & International minimum standards for bank capital adequacy, stress testing, and liquidity\\
    \hline
    \textbf{Solvency II} & EIOPA / PRA & Risk-based capital requirements and governance standards for insurers\\
    \hline
    \textbf{ICAAP Guidance} & PRA & Expectations for banks' internal capital adequacy assessment process\\
    \hline
    \textbf{ORSA Guidance} & EIOPA / PRA & Expectations for insurers' own risk and solvency assessment\\
    \hline
    \textbf{JMLSG Guidance} & JMLSG & Sector-specific guidance on AML risk assessment and controls for financial services firms\\
    \hline
    \textbf{FCA Consumer Duty} & FCA & Framework for assessing and delivering good outcomes for retail customers\\
    \hline
    \textbf{DORA} & EU & Comprehensive ICT risk management and resilience framework for EU financial services\\
    \hline
  \end{tabularx}
  \caption{Financial Services}
\end{table}

%---------------------------- SECTION ----------------------------
\section*{D.5 Climate and Sustainability}

\begin{table}[H]
  \centering
  \renewcommand{\arraystretch}{1.5}
  \begin{tabularx}{\textwidth}{ | >{\raggedright\arraybackslash}p{0.4\textwidth} 
								| >{\raggedright\arraybackslash}p{0.2\textwidth} 
                                | >{\raggedright\arraybackslash}X | }
    \hline
    \textbf{Framework / Standard} & \textbf{Published by} & \textbf{Key Purpose} \\
    \hline
    \textbf{TCFD Recommendations} & FSB / TCFD & Framework for climate-related financial risk disclosure — Governance, Strategy, Risk Management, Metrics\\
    \hline
    \textbf{TNFD Framework} & TNFD & Framework for nature-related financial risk assessment and disclosure\\
    \hline
    \textbf{European Sustainability Reporting Standards (ESRS)} & EFRAG & Mandatory sustainability reporting standards under the CSRD\\
    \hline
    \textbf{GHG Protocol} & WRI / WBCSD & Standards for measuring and reporting greenhouse gas emissions — Scopes 1, 2, and 3\\
    \hline
    \textbf{Science Based Targets initiative (SBTi)} & SBTi & Framework for setting science-aligned emissions reduction targets\\
    \hline
    \textbf{GRI Standards} & GRI & International standards for sustainability reporting\\
    \hline
  \end{tabularx}
  \caption{Climate and Sustainability}
\end{table}

%---------------------------- SECTION ----------------------------
\section*{D.6 Governance and Internal Control}

\begin{table}[H]
  \centering
  \renewcommand{\arraystretch}{1.5}
  \begin{tabularx}{\textwidth}{ | >{\raggedright\arraybackslash}p{0.4\textwidth} 
								| >{\raggedright\arraybackslash}p{0.2\textwidth} 
                                | >{\raggedright\arraybackslash}X | }
    \hline
    \textbf{Framework / Standard} & \textbf{Published by} & \textbf{Key Purpose} \\
    \hline
    \textbf{UK Corporate Governance Code} & FRC & Best practice framework for corporate governance of UK premium-listed companies\\
    \hline
    \textbf{COSO Internal Control — Integrated Framework} & COSO & Framework for designing, implementing, and assessing internal control systems\\
    \hline
    \textbf{Three Lines Model} & IIA & Framework for defining governance and risk management roles across the three lines\\
    \hline
    \textbf{IIA International Standards for Internal Auditing} & IIA & Professional standards for the practice of internal auditing\\
    \hline
  \end{tabularx}
  \caption{Governance and Internal Control}
\end{table}
